\documentclass[../main.tex]{subfiles}

\begin{document}
% -----------------------------
\section{Back end}
% -----------------------------
As a follow-up, we will create Bash scripts that automate the compilation pipeline
for our model. These scripts will take the MLIR program produced after lowering and
progressively convert it into LLVM IR, and then into a final executable. By using LLVM
tools in the scripts, we can easily compile and run our MLIR-based model from the command line.

\begin{lstlisting} [language=bash, caption={New scripts for the pipeline}]
full_compiler.sh
back_end
├── back_end.sh
\end{lstlisting}

In the \path{back_end.sh} script, implement the first step of progressive compilation.
\begin{lstlisting} [language=Python, caption={Creating LLVM IR by using mlir-opt}]
#/bin/sh
MODEL_NAME="$1"
MODEL_STEM="${MODEL_NAME%.onnx}"
echo "Compiling backend..."

# More robust way to determine script's location
SCRIPT_DIR="$(cd "$(dirname "${BASH_SOURCE[0]}")" && pwd)"
TOOLCHAIN_BIN_DIR="$SCRIPT_DIR/../../LLVM_IR/llvm-project/build/bin"
BUILD_DIR=${SCRIPT_DIR}/../build

${TOOLCHAIN_BIN_DIR}/mlir-opt ${BUILD_DIR}/${MODEL_STEM}_lowered.mlir \
  -convert-scf-to-cf \
  -convert-cf-to-llvm \
  -convert-arith-to-llvm \
  -convert-func-to-llvm \
  -reconcile-unrealized-casts \
  -o ${BUILD_DIR}/${MODEL_STEM}_llvm.mlir
echo "Created file: ${BUILD_DIR}/${MODEL_STEM}_llvm.mlir"
\end{lstlisting}
This command runs mlir-opt to apply a sequence of lowering passes that
convert the MLIR program into the LLVM dialect. The passes progressively
transform structured control flow (scf), control flow (cf), arithmetic
operations (arith), and functions (func) into their LLVM equivalents.
The result is written to a new file containing LLVM-level MLIR, which
can later be translated into LLVM IR and compiled into an executable.
Use it from our main pipeline script, together with \path{hc_main.py}.
Define \path{full_compiler.py} script with following content:
\begin{lstlisting} [language=Python, caption={Full compiler script implementation}]
#!/usr/bin/env bash

MODEL_NAME=${1:-score_model.onnx}

echo "Using model: $MODEL_NAME"

python3 hc_main.py "$MODEL_NAME"
back_end/back_end.sh "$MODEL_NAME"
\end{lstlisting}
As you can see, for back end, we are using results previously generated in \path{hc_main.py}
step, so we need to adjust that script to produce wanted files.
In \path{hc_main.py}, add imports:

\begin{lstlisting} [language=Python, caption={Needed imports}]
from hc_dialect import HiCompiler
from front_end.build_model import build_score_model, MODEL_PATH
\end{lstlisting}

Our \path{hc_main.py} now accepts model path argument, that we use to import
onnx model.
\begin{lstlisting} [language=Python, caption={Modified main function}]
def main():
    # Optional argument
    model_arg = sys.argv[1] if len(sys.argv) > 1 else MODEL_PATH
    model_path = Path(model_arg)

    # Extract base name (without extension)
    model_name = model_path.stem

    build_dir = Path("build")
    build_dir.mkdir(exist_ok=True)

    if not model_path.exists():
        candidate = build_dir / model_path.name
        if candidate.exists():
            model_path = candidate
        else:
            model_path = build_dir / model_path.name
            print(f"Model not found, building: {model_path}")
            build_score_model(str(model_path))

    original_file = build_dir / f"{model_name}_original.mlir"
    lowered_file = build_dir / f"{model_name}_lowered.mlir"

    ctx = build_context()
    module = import_onnx_to_hc_module(
        ctx, onnx_path=str(model_path), fn_name="my_func"
    )

    print("Saving high-level module...")
    with open(original_file, "w", encoding="utf-8") as f:
        f.write(str(module))

    config = MiddleEndPipelineConfig(
        apply_lowering=True,
        apply_constant_folding=False,
        apply_dce=False,
        run_analysis=False,
        debug_mode=False,
    )

    middle_end_pipeline = MiddleEndPipeline(config)
    middle_end_pipeline.apply_passes(module)

    print("Saving lowered module...")
    with open(lowered_file, "w", encoding="utf-8") as f:
        f.write(str(module))
\end{lstlisting}
Remove unnecessary imports from \path{hc_main.py}.
All generated files will be placed inside build folder.

You need to give execution rights to both newly created bash scripts:
\begin{lstlisting} [language=bash, caption={Changing execution rights}]
chmod +x full_compiler.sh
chmod +x back_end/back_end.sh
\end{lstlisting}


If you run the \path{full_compiler.sh} script now, it will fail, since we need
to fix a bug that we introduced when we implemented pow lowering. We need to
use an IndexType for \path{scf.for} upper bound.
Change lowering part of Pow operation in \path{hc_lowering} file:

\begin{lstlisting} [language=Python, caption={Final Pow lowering}]
        # ---------------- hc.pow ----------------
        if op.name == "hc.pow":
            base, exp = op.operands
            ty = op.results[0].type  # integer type

            iv_ty = builtin.IndexType()

            # loop-control constants must be index
            c0 = arith.ConstantOp.from_int_and_width(0, iv_ty)
            c1 = arith.ConstantOp.from_int_and_width(1, iv_ty)

            # exponent is i32 -> cast to index for scf.for upper bound
            exp_idx = arith.IndexCastOp(exp, iv_ty)

            # init accumulator = 1
            init = arith.ConstantOp.from_int_and_width(1, 32)

            # scf.for body arguments: (iv: index, acc: i32)
            body = Block(arg_types=[iv_ty, ty])
            iv, acc = body.args

            mul = arith.MuliOp(acc, base)
            body.add_ops([mul, scf.YieldOp(mul.result)])

            loop = scf.ForOp(
                lb=c0.result,          # lower bound
                ub=exp_idx.result,     # upper bound
                step=c1.result,        # step
                iter_args=[init.result],
                body=Region(body),
            )

            # scf.ForOp returns the iter_args results
            rewriter.replace_op(
                op,
                new_ops=[c0, c1, exp_idx, init, loop],
                new_results=[loop.results[0]],
                safe_erase=True,
            )
            return
\end{lstlisting}
Now, if we run the script, we should see something like:
\begin{lstlisting} [language=bash, caption={Output of the compiler script}]
$ ./full_compiler.sh model_score.onnx
Using model: model_score.onnx
Model not found, building: build/model_score.onnx
Saved ONNX model to: build/model_score.onnx
Saving high-level module...
Saving lowered module...
Compiling backend...
Created file: <some absolute path>/src/back_end/../build/model_score_llvm.mlir
\end{lstlisting}

Continue building our back end side. Add this to your \path{back_end.sh} script:
\begin{lstlisting} [language=bash, caption={Using mlir-translate}]
${TOOLCHAIN_BIN_DIR}/mlir-translate ${BUILD_DIR}/${MODEL_STEM}_llvm.mlir \
    -mlir-to-llvmir -o ${BUILD_DIR}/${MODEL_STEM}_out.ll
echo "Created file: ${BUILD_DIR}/${MODEL_STEM}_out.ll"
\end{lstlisting}
This command uses mlir-translate to convert the MLIR file in the LLVM dialect into LLVM IR.

\begin{itemize}
    \item The LLVM dialect is a representation used inside MLIR that models LLVM-style
operations while still staying within the MLIR infrastructure. It is designed to
look very similar to LLVM instructions but remains compatible with MLIR features
such as dialects, passes, and transformations. In other words, it is a bridge between
MLIR and LLVM, allowing MLIR programs to be gradually lowered toward LLVM.

    \item LLVM IR (Intermediate Representation) is the core low-level language used by
the LLVM compiler infrastructure. It is a platform-independent representation
that is close to machine instructions but still abstract enough to allow further
optimizations and analysis.
\end{itemize}

To produce assembly code, we can extend script with following command:

This command uses llc, the LLVM static compiler, to convert LLVM IR into assembly code for the target architecture.

\begin{lstlisting} [language=Python, caption={Using llc for assembly}]
${TOOLCHAIN_BIN_DIR}/llc -O2 -filetype=asm \
    ${BUILD_DIR}/${MODEL_STEM}_out.ll -o ${BUILD_DIR}/${MODEL_STEM}_out.s
echo "Created file: ${BUILD_DIR}/${MODEL_STEM}_out.s"
\end{lstlisting}
The input file \path{${MODEL_STEM}_out.ll} contains \texttt{LLVM IR}, and the options \texttt{-O2}
and \texttt{-filetype=asm} tell \texttt{llc} to apply level-2
optimizations and generate assembly output.

To generate an object file, add:
\begin{lstlisting} [language=Python, caption={Using llc for object file}]
${TOOLCHAIN_BIN_DIR}/llc -filetype=obj ${BUILD_DIR}/${MODEL_STEM}_out.ll \
    -o ${BUILD_DIR}/${MODEL_STEM}_out.o
echo "Created file: ${BUILD_DIR}/${MODEL_STEM}_out.o"
\end{lstlisting}
Command runs llc, the LLVM static compiler, to translate the LLVM IR file
into a target-specific object file. The option \texttt{-filetype=obj}
instructs llc to generate an object file (\texttt{.o}), which can later be
linked with other object files or libraries to produce the final executable.

Note: You can use \texttt{llvm-objdump -d \path{build/model_score_out.o}} to inspect
object file. \path{llvm-objdump} is a part of LLVM toolchain
that lets us verify that our
\path{MLIR -> LLVM -> machine code} pipeline produced the expected instructions.

\end{document}
